\documentclass{beamer}
\usepackage[utf8]{inputenc}
\usepackage[T1]{fontenc}
\usepackage[portuguese]{babel}
\usepackage{amsmath}
\usepackage{amssymb}
\usepackage{tikz}
\usetikzlibrary{shapes.geometric, arrows.meta, calc, decorations.pathreplacing, patterns}

% Configuração do Tema
\usetheme{Madrid}
\usecolortheme{default}
\beamertemplatenavigationsymbolsempty

% Comando para o círculo vermelho
\providecommand{\circled}[1]{\mbox{\tikz[baseline=(char.base)]{\node[shape=circle,draw=red,thick,inner sep=0.5pt,text=red,font=\tiny\bfseries] (char) {#1};}}}

% Convenções visuais
\tikzset{
	eixo/.style={->, thick},
	vetorv/.style={->, ultra thick, blue!75!black},
	vetora/.style={->, ultra thick, red!80!black},
	vetoraux/.style={->, thick, green!60!black},
	ponto/.style={circle, fill=black, inner sep=1.5pt}
}

% Informações do Título
\title{Aula 7: Teoria e Prática}
\subtitle{ Movimento Circular e Força centripeta}
\date{27 Abril 2026}

\begin{document}
	
	% SLIDE 1
	\frame{\titlepage}
	
	
	
	% SLIDE 2
	\begin{frame}{Resumo movimento no plano: posição e deslocamento}
		\footnotesize
		
		A localização de uma partícula no plano é especificada pelo vetor posição
		\begin{equation*}
			\vec r(t)=x(t)\vec i+y(t)\vec j.
		\end{equation*}
		
		Se a partícula passa da posição \(\vec r_1\) para a posição \(\vec r_2\), o deslocamento é
		\begin{equation*}
			\Delta \vec r=\vec r_2-\vec r_1=\Delta x\,\vec i+\Delta y\,\vec j.
		\end{equation*}
		
		\textcolor{red}{\textbf{Ideia importante:}} embora o movimento seja vetorial, podemos estudar separadamente as componentes \(x\) e \(y\), desde que no final interpretemos o resultado como vetor.
		
		\vspace{0.15cm}
		\begin{columns}[c]
			\begin{column}{0.55\textwidth}
				\textcolor{blue}{\textbf{Leitura geométrica:}}
				\begin{itemize}
					\item \(x(t)\) mede a posição ao longo do eixo \(x\);
					\item \(y(t)\) mede a posição ao longo do eixo \(y\);
					\item \(\Delta \vec r\) liga diretamente a posição inicial à posição final;
					\item a trajetória pode ser curva, mas \(\Delta \vec r\) é sempre um vetor entre dois pontos.
				\end{itemize}
			\end{column}
			\begin{column}{0.45\textwidth}
				\centering
				\begin{tikzpicture}[scale=0.75]
					\draw[eixo] (-0.2,0) -- (4.2,0) node[right] {$x$};
					\draw[eixo] (0,-0.2) -- (0,3.4) node[above] {$y$};
					\coordinate (A) at (0.9,0.7);
					\coordinate (B) at (3.4,2.6);
					\draw[thick, blue!60!black] plot[smooth] coordinates {(A) (1.4,1.5) (2.0,1.2) (2.7,2.2) (B)};
					\filldraw[ponto] (A) node[below left] {$A$};
					\filldraw[ponto] (B) node[above right] {$B$};
					\draw[vetora] (A) -- (B) node[midway, below right] {$\Delta\vec r$};
					\node[font=\scriptsize, blue!70!black] at (2.7,0.8) {trajetória};
				\end{tikzpicture}
			\end{column}
		\end{columns}
		
	\end{frame}
	
	
	% SLIDE 3
	\begin{frame}{Exemplo 4.01 (Hal-Res): vetor posição do coelho}
		\scriptsize
		
		Um coelho atravessa um estacionamento. As coordenadas da posição são dadas, no SI, por
		\begin{equation*}
			x(t)= -0{,}31t^2+7{,}2t+28, \qquad
			y(t)= 0{,}22t^2-9{,}1t+30.
		\end{equation*}
		Logo,
		\begin{equation*}
			\vec r(t)=\left(-0{,}31t^2+7{,}2t+28\right)\vec i+
			\left(0{,}22t^2-9{,}1t+30\right)\vec j.
		\end{equation*}
		
		\begin{columns}[T]
			\begin{column}{0.55\textwidth}
				\textcolor{blue}{\textbf{No instante \(t=15\,\text{s}\):}}
				\begin{align*}
					x(15)&=-0{,}31(15)^2+7{,}2(15)+28 \approx 66\,\text{m},\\
					y(15)&=0{,}22(15)^2-9{,}1(15)+30=-57\,\text{m}.
				\end{align*}
				Portanto,
				\begin{equation*}
					\textcolor{red}{\vec r(15\,\text{s})=(66\,\text{m})\vec i-(57\,\text{m})\vec j.}
				\end{equation*}
				O módulo e o ângulo são:
				\begin{equation*}
					|\vec r|=\sqrt{(66)^2+(-57)^2}\approx 87\,\text{m},
				\end{equation*}
				\begin{equation*}
					\theta=\arctan\left(\frac{-57}{66}\right)\approx -41^\circ.
				\end{equation*}
			\end{column}
			\begin{column}{0.45\textwidth}
				\centering
				\begin{tikzpicture}[scale=0.027]
					\draw[->, thick] (-20,0) -- (125,0) node[right] {$x$};
					\draw[->, thick] (0,-105) -- (0,60) node[above] {$y$};
					
					\draw[domain=0:24, smooth, variable=\t, thick, blue!70!black]
					plot ({-0.31*\t*\t+7.2*\t+28},{0.22*\t*\t-9.1*\t+30});
					
					\coordinate (O) at (0,0);
					\coordinate (Px) at (66.25,0);
					\coordinate (P) at (66.25,-57);
					
					% vetor posição
					\draw[->, thick, red] (O) -- (P) node[midway, below left] {$\vec r(15)$};
					
					% projeções nas componentes
					\draw[dashed, gray] (P) -- (Px);
					
					
					% ponto
					\filldraw[black] (P) circle (3pt) node[right] {$t=15\,\text{s}$};
					
					% arco do ângulo
					\draw[->, thick, black] (18,0) arc[start angle=0,end angle=-41,radius=18];
					\node[font=\tiny] at (37,-8) {$\theta=-41^\circ$};
					
					\node[font=\scriptsize] at (72,25) {trajetória do coelho};
				\end{tikzpicture}
				
				\vspace{0.1cm}
				\textcolor{red}{\textbf{Interpretação:}} o coelho está no quarto quadrante: \(x>0\) e \(y<0\).
			\end{column}
		\end{columns}
		
	\end{frame}
	
	% SLIDE 4
	\begin{frame}{Exemplo 4.02(Hal-Res): velocidade do coelho}
		\scriptsize
		
		A velocidade instantânea é a derivada do vetor posição:
		$\vec v(t)
		=\frac{dx}{dt}\vec i+\frac{dy}{dt}\vec j
		=v_x(t)\vec i+v_y(t)\vec j.$
		
		\begin{equation*}
			v_x(t)=\frac{d}{dt}\left(-0{,}31t^2+7{,}2t+28\right)=-0{,}62t+7{,}2,
			\qquad
			v_y(t)=\frac{d}{dt}\left(0{,}22t^2-9{,}1t+30\right)=0{,}44t-9{,}1.
		\end{equation*}
		
		\begin{columns}[T]
			\begin{column}{0.70\textwidth}
				\textcolor{blue}{\textbf{No instante \(t=15\,\text{s}\):}}
				\begin{equation*}
					v_x=-2{,}1\,\text{m/s}, \qquad v_y=-2{,}5\,\text{m/s}.
				\end{equation*}
				\begin{equation*}
					\textcolor{red}{\vec v(15\,\text{s})=(-2{,}1\,\text{m/s})\vec i-(2{,}5\,\text{m/s})\vec j.}
				\end{equation*}
				
				O módulo é:				$v=\sqrt{v_x^2+v_y^2}
				=\sqrt{(-2{,}1)^2+(-2{,}5)^2}
				\approx 3{,}3\,\text{m/s}.
				$
				
				O ângulo satisfaz:
				\[
				\theta=\arctan\left(\frac{v_y}{v_x}\right)
				=\arctan\left(\frac{-2{,}5}{-2{,}1}\right)=-130^\circ.
				\]
				
				Como \(v_x<0\) e \(v_y<0\), o vetor está no \textbf{terceiro quadrante}.  	Logo, \textcolor{red}{\(\theta\approx -130^\circ\)}.\textcolor{red}{\textbf{Interpretação:}} a velocidade é tangente à trajetória no instante considerado e aponta para a esquerda e para baixo.
			\end{column}
			
			\begin{column}{0.35\textwidth}
				\centering
				\begin{tikzpicture}[x=0.028cm,y=0.028cm,>=stealth,font=\scriptsize]
					% Eixos
					\draw[->, thick] (0,0) -- (88,0) node[right] {$x\;(\text{m})$};
					\draw[->, thick] (0,-72) -- (0,45) node[above] {$y\;(\text{m})$};
					
					% Trajetória suave
					\draw[thick, cyan!70!black, domain=0:25, samples=120, smooth, variable=\t]
					plot ({-0.31*\t*\t+7.2*\t+28},{0.22*\t*\t-9.1*\t+30});
					
					% Ponto no instante t=15 s
					\coordinate (P) at (66.25,-57);
					\filldraw[orange!85!brown] (P) circle (3pt);
					\node[right=11pt] at (P) {$t=15\,\text{s}$};
					
					% Vetor velocidade
					\coordinate (Vtip) at (49.5,-77);
					\draw[->, thick, magenta!80!black] (P) -- (Vtip)
					node[below left] {$\vec v$};
					
					% Componentes
					\coordinate (VxEnd) at (49.5,-57);
					\coordinate (VyEnd) at (66.25,-77);
					
					\draw[->, thick, magenta!60!black] (P) -- (VxEnd);
					\draw[->, thick, magenta!60!black] (P) -- (VyEnd);
					
					% Só a linha tracejada vertical auxiliar
					\draw[dashed, gray] (Vtip) -- (VxEnd);
					
					% Linha cinzenta de referência do ângulo
					% começa um pouco à direita de P para não cortar v_y
					\draw[thin, gray!80] ($(P)+(2.5,0)$) -- ($(P)+(15,0)$);
					
					% Arco do ângulo
					\draw[thin, gray!80!black] ($(P)+(8,0)$)
					arc[start angle=0,end angle=-130,radius=8];
					\node[gray!80!black] at (82,-68.5) {$-130^\circ$};
					
					% Legenda
					\node at (53,28) {trajetória do coelho};
				\end{tikzpicture}
			\end{column}
		\end{columns}
		
		
		
	\end{frame}
	
	
	% SLIDE 5
	\begin{frame}{Exemplo 4.03 (Hal-Res): aceleração do coelho}
		\scriptsize
		
		A aceleração instantânea é a derivada da velocidade:
		\[
		\vec a(t)=\frac{d\vec v}{dt}
		=\frac{dv_x}{dt}\vec i+\frac{dv_y}{dt}\vec j
		=a_x\vec i+a_y\vec j.
		\]
		
		Como
		\[
		v_x(t)=-0{,}62t+7{,}2,
		\qquad
		v_y(t)=0{,}44t-9{,}1,
		\]
		obtemos
		\[
		a_x=\frac{dv_x}{dt}=-0{,}62\,\text{m/s}^2,
		\qquad
		a_y=\frac{dv_y}{dt}=0{,}44\,\text{m/s}^2.
		\]
		
		\begin{columns}[T]
			\begin{column}{0.64\textwidth}
				
				Logo,
				\[
				\textcolor{red}{\vec a=(-0{,}62\,\text{m/s}^2)\vec i+(0{,}44\,\text{m/s}^2)\vec j.}
				\]
				
				O módulo é
				\[
				a=\sqrt{a_x^2+a_y^2}
				=\sqrt{(-0{,}62)^2+(0{,}44)^2}
				\approx 0{,}76\,\text{m/s}^2.
				\]
				
				O ângulo calculado diretamente é
				\[
				\theta=\arctan\left(\frac{a_y}{a_x}\right)
				=\arctan\left(\frac{0{,}44}{-0{,}62}\right)
				\approx -35^\circ.
				\]
				
				\vspace{0.1cm}
				\textcolor{red}{\textbf{Interpretação:}} a aceleração é constante e indica como o vetor velocidade está a mudar.
				
			\end{column}
			
			\begin{column}{0.36\textwidth}
				\centering
				\begin{tikzpicture}[x=0.025cm,y=0.025cm,>=stealth,font=\scriptsize]
					% Eixos
					\draw[->, thick] (0,0) -- (88,0) node[right] {$x\;(\text{m})$};
					\draw[->, thick] (0,-72) -- (0,45) node[above] {$y\;(\text{m})$};
					
					% Trajetória suave
					\draw[thick, cyan!70!black, domain=0:25, samples=140, smooth, variable=\t]
					plot ({-0.31*\t*\t+7.2*\t+28},{0.22*\t*\t-9.1*\t+30});
					
					% Ponto t=15 s
					\coordinate (P) at (66.25,-57);
					\filldraw[orange!85!brown] (P) circle (3pt);
					\node[right=3pt] at (P) {$t=15\,\text{s}$};
					
					% Vetor aceleração
					\coordinate (Atip) at (50,-45);
					\draw[->, semithick, orange!85!black] (P) -- (Atip)
					node[above left] {$\vec a$};
					
					% Componentes
					\coordinate (AxEnd) at (50,-57);
					\coordinate (AyEnd) at (66.25,-45);
					
					\draw[->, semithick, orange!70!black] (P) -- (AxEnd);
					\draw[->, semithick, orange!70!black] (P) -- (AyEnd);
					
					% Retângulo auxiliar
					\draw[dashed, gray] (Atip) -- (AxEnd);
					\draw[dashed, gray] (Atip) -- (AyEnd);
					
					% Linha horizontal de referência para o ângulo, no sentido de a_x
					\draw[thin, gray!70] (P) -- ++(-18,0);
					
					% Arco do ângulo -35° entre a direção horizontal para a esquerda e o vetor a
					\draw[thin, green!80!black]
					($(P)+(-13,0)$)
					arc[start angle=180, delta angle=-35, radius=13];
					
					\node[green!80!black] at ($(P)+(-30,5)$) {$-35^\circ$};
					
					
					
					% Legenda
					\node at (53,28) {trajetória do coelho};
				\end{tikzpicture}
			\end{column}
		\end{columns}
		
	\end{frame}
	
	
	
	% SLIDE 6
	\begin{frame}{Movimento circular uniforme: definição}
		\scriptsize
		
		Uma partícula está em \textbf{movimento circular uniforme} quando descreve uma circunferência, ou um arco de circunferência, com velocidade escalar constante. \textcolor{red}{\textbf{Ponto essencial:}} mesmo com velocidade escalar constante, a partícula está acelerada porque a direção da velocidade muda continuamente.
		
		\begin{columns}[c]
			\begin{column}{0.65\textwidth}
				\begin{itemize}
					\item \(\vec v\): tangente à trajetória e \(|\vec v|=v\) constante;
					\item \(\vec a_c\): radial, apontando para o centro, e \(|\vec a_c|\) constante.
				\end{itemize}
				
				\begin{equation*}
					\vec a_{\text{méd}}
					=
					\frac{\Delta \vec v}{\Delta t}
					=
					\frac{\vec v_2-\vec v_1}{\Delta t}.
				\end{equation*}
				
				Como \(\vec v_1\) e \(\vec v_2\) têm o mesmo módulo, mas direções diferentes (a diferença é vetorial):$\Delta \vec v=\vec v_2-\vec v_1\neq 0.$
				
				\vspace{0.2cm}
				
				Logo, \textcolor{red}{\(\vec a\neq 0\)}. No limite \(\Delta t\to 0\), \(\Delta\vec v\) aponta para o centro; por isso a \textbf{aceleração é centrípeta}.
			\end{column}
			
			\begin{column}{0.35\textwidth}
				\centering
				\begin{tikzpicture}[scale=0.6,>=stealth]
					
					% centro e raio
					\coordinate (C) at (0,0);
					\def\R{2.0}
					
					% circunferência
					\draw[thin, orange!55!brown] (C) circle (\R);
					
					% setas de sentido ao longo da circunferência
					% (anti-horário e afastadas dos pontos materiais)
					\draw[->, thick, orange!55!brown]
					(55:\R) arc[start angle=55,end angle=80,radius=\R];
					
					\draw[->, thick, orange!55!brown]
					(170:\R) arc[start angle=170,end angle=195,radius=\R];
					
					\draw[->, thick, orange!55!brown]
					(295:\R) arc[start angle=295,end angle=320,radius=\R];
					
					% posições da partícula
					\coordinate (P1) at (135:\R);
					\coordinate (P2) at (0:\R);
					\coordinate (P3) at (230:\R);
					
					% centro
					\filldraw[black] (C) circle (1.8pt);
					\node[below left=1pt] at (C) {\scriptsize centro};
					
					% linhas radiais tracejadas
					\draw[dashed, gray!70] (P1) -- (C);
					\draw[dashed, gray!70] (P2) -- (C);
					\draw[dashed, gray!70] (P3) -- (C);
					
					% partículas
					\shade[ball color=cyan!55, draw=cyan!35!white] (P1) circle (0.18);
					\shade[ball color=cyan!55, draw=cyan!35!white] (P2) circle (0.18);
					\shade[ball color=cyan!55, draw=cyan!35!white] (P3) circle (0.18);
					
					% vetores velocidade
					\draw[->, magenta!85!black] (P1) -- +(-1.35,-0.85)
					node[left] {\(\vec v\)};
					
					\draw[->, magenta!85!black] (P2) -- +(0,1.55)
					node[above] {\(\vec v\)};
					
					\draw[->, magenta!85!black] (P3) -- +(1.25,-0.85)
					node[right] {\(\vec v\)};
					
					% vetores aceleração centrípeta
					\draw[->, orange!90!black] (P1) -- ($(P1)!0.42!(C)$)
					node[midway, right] {\(\vec a\)};
					
					\draw[->, thick, orange!90!black] (P2) -- ($(P2)!0.42!(C)$)
					node[midway, below] {\(\vec a\)};
					
					\draw[->, orange!90!black] (P3) -- ($(P3)!0.42!(C)$)
					node[midway, left] {\(\vec a\)};
					
				\end{tikzpicture}
			\end{column}
		\end{columns}
		
		\vspace{0.2cm}
		
		O módulo é	$\textcolor{red}{a_c=\frac{v^2}{r}}$. O tempo para completar uma volta é o \textbf{período} \(T\). Como a distância de uma volta completa é \(2\pi r\), temos $\textcolor{red}{T= \frac{\text{distância de uma volta completa}}{v}=\frac{2\pi r}{v}}.$
		
		\vspace{0.2cm}
		\textcolor{blue}{\textbf{Exemplo simples:}} Um velocista corre com \(v=10\,\text{m/s}\) e faz uma curva de raio \(r=25\,\text{m}\). Qual é o módulo da aceleração centrípeta? Resposta: $	a_c=\frac{v^2}{r}=\frac{(10)^2}{25}=4{,}0\,\text{m/s}^2.$
		
		\vspace{0.2cm}
		
		\textcolor{red}{\textbf{Não confundir:}} \(v\) constante \(\not\Rightarrow \vec v\) constante. O módulo da velocidade é constante, mas o vetor velocidade muda de direção.
		
	\end{frame}
	
	% SLIDE 7
	\begin{frame}{Força centrípeta}
		\scriptsize
		
		Pela Segunda Lei de Newton, se existe uma aceleração radial para o centro, deve existir uma força resultante \(F_{\text{res}}\) na mesma direção e no mesmo sentido:
		\[
		\vec F_{\text{res}}=m\vec a_c,
		\qquad
		\textcolor{red}{F_c=m\frac{v^2}{R}}.
		\]
		
		\textcolor{red}{\textbf{Importante:}} a ``força centrípeta'' não é uma nova força da natureza. É o nome dado à \textbf{resultante radial} das forças reais \textbf{para realizar o movimento circular}.
		
		\vspace{0.1cm}
		\begin{columns}[T,onlytextwidth]
			
			%====================================================
			\begin{column}{0.49\textwidth}
				\textcolor{blue}{\textbf{Exemplo 1: carro numa curva plana}}
				
				\vspace{0.05cm}
				\begin{columns}[c,onlytextwidth]
					\begin{column}{0.42\textwidth}
						\tiny
						O carro precisa de uma força radial para fazer a curva. Essa força é fornecida pelo \textbf{atrito} entre os pneus e a estrada.
						
						\vspace{0.15cm}
						\[
						f=m\frac{v^2}{R}
						\]
						
						\vspace{-0.15cm}
						\begin{center}
							O atrito fornece\\
							a resultante radial.
						\end{center}
					\end{column}
					\begin{column}{0.58\textwidth}
						\centering
						\begin{tikzpicture}[scale=0.6,>=stealth]
							% centro e raio da curva
							\coordinate (C) at (0,0);
							\def\R{2.2}
							
							% estrada (arco)
							\draw[thick, gray!65] (205:\R) arc[start angle=205,end angle=100,radius=\R];
							
							% posição do carro
							\coordinate (Car) at (225:\R);
							
							% centro
							\filldraw[black] (C) circle (1.2pt) node[right] {\tiny centro};
							\draw[dashed, gray] (Car) -- (C);
							
							% carro
							\draw[fill=blue!20, rotate around={45:(Car)}]
							($(Car)+(-0.35,-0.18)$) rectangle ($(Car)+(0.35,0.18)$);
							\node[font=\tiny] at (Car) {carro};
							
							% velocidade tangente (movimento horário)
							\draw[->, thick, magenta!85!black]
							(Car) -- +(135:1.0)
							node[above left=-1pt] {$\vec v$};
							
							% força de atrito para o centro
							\draw[->, thick, orange!90!black]
							(Car) -- ($(Car)!0.55!(C)$)
							node[midway, above] {$\vec f$};
						\end{tikzpicture}
					\end{column}
				\end{columns}
			\end{column}
			
			%====================================================
			\begin{column}{0.49\textwidth}
				\textcolor{blue}{\textbf{Exemplo 2: pedra presa a uma corda}}
				
				\vspace{0.05cm}
				\begin{columns}[c,onlytextwidth]
					\begin{column}{0.42\textwidth}
						\tiny
						A pedra move-se em circunferência porque a corda exerce uma \textbf{tração} dirigida para o centro.
						
						\vspace{0.15cm}
						\[
						T=m\frac{v^2}{R}
						\]
						
						\vspace{-0.15cm}
						\begin{center}
							A tração fornece\\
							a resultante radial.
						\end{center}
					\end{column}
					\begin{column}{0.58\textwidth}
						\centering
						\begin{tikzpicture}[scale=0.6,>=stealth]
							\coordinate (O) at (0,0);
							\def\R{1.6}
							\coordinate (P) at (40:\R);
							
							% trajetória
							\draw[thin, orange!55!brown] (O) circle (\R);
							\draw[->, thick, orange!55!brown]
							(255:\R) arc[start angle=255,end angle=295,radius=\R];
							
							% mão e corda
							\filldraw[black] (O) circle (1.3pt) node[below] {\tiny mão};
							\draw[thick] (O) -- (P);
							
							% pedra
							\shade[ball color=gray!35, draw=black] (P) circle (0.16);
							\node[right=2pt] at (P) {\tiny pedra};
							
							% velocidade e tração
							\draw[->, thick, magenta!85!black]
							(P) -- +(130:0.8)
							node[above] {$\vec v$};
							
							\draw[->, thick, orange!90!black]
							(P) -- ($(P)!0.48!(O)$)
							node[midway, below] {$\vec T$};
						\end{tikzpicture}
					\end{column}
				\end{columns}
			\end{column}
			
		\end{columns}
		
		\vspace{0.05cm}
		\textcolor{red}{No \textbf{DCL:}} desenhamos as forças reais, como \(\vec f\) ou \(\vec T\), e depois identificamos a componente radial da resultante: $		\sum F_r=m\frac{v^2}{R}.$
		
	\end{frame}
	
	
	% SLIDE 8
	\begin{frame}{Derivação qualitativa da aceleração centrípeta}
		\scriptsize
		
		Considere uma partícula em \textbf{movimento circular uniforme}, de raio \(r\) e rapidez constante \(v\). O módulo da velocidade não muda, mas a sua \textbf{direção} muda continuamente.
		
		\vspace{0.05cm}
		\begin{columns}[T]
			\begin{column}{0.58\textwidth}
				
				Num pequeno intervalo \(\Delta t\), a partícula percorre um arco \(\Delta s\) e roda de um ângulo \(\Delta\theta\). Os vetores \(\vec v_1\) e \(\vec v_2\) têm o mesmo módulo \(v\), mas direções diferentes.
				
				\vspace{0.3cm}
				Para ângulos pequenos:
				$
				|\Delta \vec v|\approx v\,\Delta\theta.
				$
				\vspace{0.1cm}
				A aceleração média é:
				$
				a_c\approx \frac{|\Delta \vec v|}{\Delta t}
				\approx
				\frac{v\,\Delta\theta}{\Delta t}.
				$
				
				Como o arco percorrido é
				$
				\Delta s=r\,\Delta\theta,
				$
				então, dividindo por \(\Delta t\):
				\[
				\frac{\Delta s}{\Delta t}
				=
				r\frac{\Delta\theta}{\Delta t}.
				\]
				
				Mas \(\frac{\Delta s}{\Delta t}=v\). Portanto:
				$
				\frac{\Delta\theta}{\Delta t}=\frac{v}{r}.
				$
				
				Logo:
				\[
				\textcolor{red}{a_c=v\frac{v}{r}=\frac{v^2}{r}.}
				\]
				
				\vspace{0.05cm}
				\textcolor{red}{\textbf{Interpretação:}} \(\Delta\vec v\) aponta para o interior da trajetória; no limite \(\Delta t\to 0\), a aceleração aponta para o centro.
				
			\end{column}
			
			\begin{column}{0.42\textwidth}
				\centering
				\begin{tikzpicture}[scale=0.72,>=stealth]
					\coordinate (C) at (0,0);
					\def\R{1.75}
					
					% circunferência
					\draw[thin, orange!55!brown] (C) circle (\R);
					\filldraw[ponto] (C) node[below left] {\scriptsize centro};
					
					% pontos próximos
					\coordinate (P1) at (25:\R);
					\coordinate (P2) at (55:\R);
					
					% raios
					\draw[dashed, gray] (C) -- (P1);
					\draw[dashed, gray] (C) -- (P2);
					
					% partículas
					\filldraw[blue!60!cyan] (P1) circle (2pt);
					\filldraw[blue!60!cyan] (P2) circle (2pt);
					
					% velocidades
					\draw[->, thick, magenta!85!black]
					(P1) -- +({-0.8*sin(25)},{0.8*cos(25)})
					node[right] {\(\vec v_1\)};
					
					\draw[->, thick, magenta!85!black]
					(P2) -- +({-0.8*sin(55)},{0.8*cos(55)})
					node[above left] {\(\vec v_2\)};
					
					% arco angular
					\draw[->, thick, gray!80!black] (0.65,0) arc[start angle=0,end angle=25,radius=0.65];
					\node at (0.95,0.18) {\(\Delta\theta\)};
					
					% arco percorrido
					\draw[ultra thick, red!80!black] (25:\R) arc[start angle=25,end angle=55,radius=\R];
					\node[red!80!black] at (1.55,1.15) {\(\Delta s\)};
					
					% triângulo auxiliar de velocidades (rodado de 90 graus)
					\begin{scope}[shift={(0.95,-3.55)}, scale=0.9, rotate=90]
						\coordinate (A) at (0,0);
						\coordinate (B) at (1.2,0);
						\coordinate (D) at (0.95,0.58);
						
						\draw[->, thick, magenta!85!black] (A) -- (B) node[midway, right] {\(\vec v_1\)};
						\draw[->, thick, magenta!85!black] (A) -- (D) node[midway, left] {\(\vec v_2\)};
						\draw[->, thick, orange!90!black] (B) -- (D) node[midway, above right] {\(\Delta\vec v\)};
						
						\node[font=\scriptsize] at (-0.5,-2.55) {triângulo das velocidades};
					\end{scope}
					
				\end{tikzpicture}
			\end{column}
		\end{columns}
		
	\end{frame}
	
	% SLIDE 9
	\begin{frame}{Movimento circular uniforme e uniformemente variado}
		\scriptsize
		
		Considere uma \textbf{partícula} que se move sobre uma circunferência de raio \(r\).  
		Em vez de descrever a posição pelo comprimento do arco \(s\), podemos descrever a posição pela \textbf{posição angular} \(\theta\). ($\textcolor{red}{1\,\text{volta}=2\pi\,\text{rad}=360^\circ}$)
		
		\vspace{0.1cm}
		\begin{columns}[c]
			\begin{column}{0.58\textwidth}
				\textcolor{blue}{\textbf{Descrição escalar ao longo da trajetória:}}
				\[
				\theta=\frac{s}{r}, 
				\qquad
				\frac{\theta}{2\pi}=\frac{\theta_{\text{graus}}}{360^\circ},
				\qquad
				\Delta\theta=\theta_2-\theta_1=\frac{\Delta s}{r}.
				\]
				
				\vspace{0.1cm}
				\textcolor{blue}{\textbf{Descrição vetorial no plano:}}
				\[
				\vec r(t)=x(t)\vec i+y(t)\vec j.
				\]
				
				Para uma trajetória circular de raio \(r\), centrada na origem:
				\[
				x(t)=r\cos\theta(t),
				\qquad
				y(t)=r\sin\theta(t).
				\]
				
				\vspace{0.05cm}
				\textcolor{red}{\textbf{Ideia-chave:}} a descrição escalar \(s\leftrightarrow\theta\) é equivalente à descrição vetorial \(\vec r(t)\), mas a forma vetorial mostra explicitamente as coordenadas no plano.
			\end{column}
			
			\begin{column}{0.42\textwidth}
				\centering
				\begin{tikzpicture}[scale=0.85]
					\coordinate (C) at (0,0);
					\def\R{1.8}
					\draw[thick] (C) circle (\R);
					\filldraw[ponto] (C) node[below left] {centro};
					
					\coordinate (A) at (\R,0);
					\coordinate (P) at (55:\R);
					\coordinate (Px) at ({\R*cos(55)},0);
					\coordinate (Py) at (0,{\R*sin(55)});
					
					% eixos
					\draw[->, thick] (-0.3,0) -- (2.4,0) node[right] {$x$};
					\draw[->, thick] (0,-0.3) -- (0,2.3) node[above] {$y$};
					
					% raio e arco
					\draw[thick] (C) -- (A);
					\draw[thick, blue!70!black] (C) -- (P) node[midway, above left] {$r$};
					
					\draw[ultra thick, red!80!black] (A) arc[start angle=0,end angle=55,radius=\R];
					\node[red!80!black] at (1.45,0.75) {$s$};
					
					\draw[->, thick, gray!80!black] (0.6,0) arc[start angle=0,end angle=55,radius=0.6];
					\node at (0.9,0.28) {$\theta$};
					
					% projeções
					\draw[dashed, gray] (P) -- (Px) node[below] {$x$};
					\draw[dashed, gray] (P) -- (Py) node[left] {$y$};
					
					% ponto
					\filldraw[blue!70!black] (P) circle (1.8pt);
					\node[above right] at (P) {partícula};
					
					% vetor posição
					\draw[->, thick, red!80!black] (C) -- (P);
				\end{tikzpicture}
			\end{column}
		\end{columns}
		
	\end{frame}
	
	
	% SLIDE 10
	\begin{frame}{}
		\scriptsize
		
	As grandezas angulares são obtidas por derivação de $\theta(t)$, analogamente ao movimento retilíneo. Indicamos com $	\omega$ a velocidade angular e com $\alpha$ a aceleração angular.
		
		\vspace{0.1cm}
		\begin{columns}[T]
			\begin{column}{0.48\textwidth}
				\textcolor{blue}{\textbf{Grandezas angulares:}}
				\[
				\omega_{\text{méd}}=\frac{\Delta\theta}{\Delta t},
				\qquad
				\omega=\frac{d\theta}{dt}.
				\]
				
				\[
				\alpha_{\text{méd}}=\frac{\Delta\omega}{\Delta t},
				\qquad
				\alpha=\frac{d\omega}{dt}
				=
				\frac{d^2\theta}{dt^2}.
				\]
				
				\vspace{0.1cm}
				\textcolor{blue}{\textbf{Analogia com a cinemática linear:}}
				\[
				x \longleftrightarrow \theta,
				\qquad
				v \longleftrightarrow \omega,
				\qquad
				a \longleftrightarrow \alpha.
				\]
				
				\textcolor{red}{\textbf{Unidades:}} 
				\(\theta\) em rad, \(\omega\) em rad/s, \(\alpha\) em rad/s\(^2\).
				
				\vspace{0.1cm}
				\textcolor{red}{\textbf{Sentido positivo:}} também no movimento circular temos um sentido positivo e um sentido negativo, pois podemos girar nos dois sentidos. Usualmente, o sentido anti-horário é escolhido como positivo.
			\end{column}
			
			\begin{column}{0.52\textwidth}
				\textcolor{blue}{\textbf{Da forma angular à forma vetorial:}}
				
				Para \(\vec r(t)=r\cos\theta(t)\vec i+r\sin\theta(t)\vec j\), a velocidade é
				\[
				\vec v(t)=\frac{d\vec r}{dt}.
				\]
				
				Como \(\omega=d\theta/dt\),
				\[
				\textcolor{red}{
					\vec v(t)=
					r\omega\left[-\sin\theta(t)\vec i+\cos\theta(t)\vec j\right].}
				\]
				
				Assim, o módulo da velocidade é
				\[
				v=\omega r \rightarrow F_c= m \omega^2 r
				\]
				
				\vspace{0.05cm}
				\textcolor{blue}{\textbf{Aceleração:}}
				$\vec a(t)=\frac{d\vec v}{dt}
				=
				\vec a_t+\vec a_r.
				$
				
				\[
				\textcolor{red}{a_t= r \frac{d \omega}{dt}=r\alpha},
				\qquad
				\textcolor{red}{a_r=\omega^2r=\frac{v^2}{r}}.
				\]
				
				\textcolor{red}{\textbf{Leitura física:}}  
				\(\vec a_t\) aceleração tangetial muda o módulo de \(\vec v\);  
				\(\vec a_r\) aceleração radial muda a direção de \(\vec v\).
			\end{column}
		\end{columns}
		
		\vspace{0.2cm}
		\centering
		Se \(\omega\) é constante: \textbf{MCU}. \qquad
		Se \(\alpha\) é constante: \textbf{MCUV}.
		
	\end{frame}
	
	
	% SLIDE 11
	\begin{frame}{Movimento circular uniforme (MCU)}
		\scriptsize
		
		No \textbf{movimento circular uniforme}, a rapidez é constante e a velocidade angular também:
		\[
		\omega=\text{constante}, \qquad \alpha=0.
		\]
		
		\begin{columns}[T,onlytextwidth]
			\begin{column}{0.56\textwidth}
				\textcolor{blue}{\textbf{Relações cinemáticas:}}
				\begin{align*}
					v &= \omega r, \qquad
					T = \frac{2\pi}{\omega}=\frac{2\pi r}{v} \rightarrow
					f = \frac{1}{T},\, \omega= \frac{2\pi}{T}= 2\pi f\\[0.1cm]
					&\text{\(f\) é a frequência: número de voltas por segundo.}\\ 
					&\text{\(T\) é o intervalo de tempo para cumprir uma volta.} 
				\end{align*}
				
				\textcolor{blue}{\textbf{Aceleração centrípeta:}}
				$
				a_c=\frac{v^2}{r}=\omega^2 r.
				$
				
				\vspace{0.2cm}
				
				\textcolor{blue}{\textbf{Dinâmica:}}
				$
				\sum F_r = m a_c = m\frac{v^2}{r}=m\omega^2 r.
				$
				
				\vspace{0.2cm}
				
				\textcolor{red}{\textbf{Exemplo simples:}} se uma partícula percorre uma circunferência de raio \(2,0\,\text{m}\) com \(v=4,0\,\text{m/s}\),
				\[
				a_c=\frac{(4,0)^2}{2,0}=8,0\,\text{m/s}^2.
				\]
			\end{column}
			
			\begin{column}{0.44\textwidth}
				\centering
				\begin{tikzpicture}[scale=0.85,>=stealth]
					\coordinate (C) at (0,0);
					\def\R{1.8}
					
					\draw[thick, orange!60!brown] (C) circle (\R);
					\filldraw[ponto] (C) node[below left] {centro};
					
					\coordinate (P) at (25:\R);
					\filldraw[blue!60!cyan] (P) circle (2.3pt);
					
					\draw[dashed, gray] (C) -- (P) node[midway, below] {$r$};
					
					% velocidade tangente à trajetória
					\draw[->, thick, magenta!85!black]
					(P) -- +({-0.95*sin(25)},{0.95*cos(25)})
					node[above] {$\vec v$};
					
					% aceleração centrípeta
					\draw[->, thick, orange!90!black]
					(P) -- ($(P)!0.5!(C)$)
					node[midway, below] {$\vec a_c$};
					
					% sentido percorrido na circunferência
					\draw[->, thick, orange!60!brown]
					(110:\R) arc[start angle=110,end angle=140,radius=\R];
					
					\node[align=center, font=\scriptsize] at (0,-2.45)
					{\(\vec v\) é tangente\\ \(\vec a_c\) aponta para o centro};
				\end{tikzpicture}
			\end{column}
		\end{columns}
		
	\end{frame}
	
	% SLIDE 12
	\begin{frame}{Movimento circular uniformemente variado (MCUV)}
		\scriptsize
		
		Quando a \textbf{aceleração angular} é constante,
		\[
		\alpha=\text{constante},
		\]
		as equações são análogas às do movimento retilíneo uniformemente variado.
		
		\vspace{0.1cm}
		\begin{columns}[T,onlytextwidth]
			\begin{column}{0.54\textwidth}
				\textcolor{blue}{\textbf{Equações angulares:}}
				\begin{align*}
					\omega(t) &= \omega(0)+\alpha t,\,\,
					\theta(t) = \theta(0)+\omega(0)t+\frac{1}{2}\alpha t^2,\\
					\omega^2(t) &= \omega^2(0)+2\alpha\bigl(\theta(t)-\theta(0)\bigr),\\
					\theta(t) &= \theta(0)+\frac{1}{2}\bigl(\omega(0)+\omega(t)\bigr)t,\,\,
					\omega_{\text{méd}} = \frac{1}{2}\bigl(\omega(0)+\omega(t)\bigr).
				\end{align*}
				
				\textcolor{blue}{\textbf{Acelerações do ponto:}}
				\begin{align*}
					a_t &= \alpha r \qquad \text{(tangencial)},\\
					a_r &= \omega^2 r \qquad \text{(radial)}.
				\end{align*}
				
				\textcolor{red}{\textbf{Interpretação:}}  
				\(a_t\) muda o \textbf{módulo} de \(\vec v\);  
				\(a_r\) muda a \textbf{direção} de \(\vec v\).
			\end{column}
			
			\begin{column}{0.46\textwidth}
				\centering
				\begin{tikzpicture}[scale=0.8,>=stealth]
					\coordinate (C) at (0,0);
					\def\R{1.7}
					\draw[thick] (C) circle (\R);
					\filldraw[ponto] (C);
					
					\coordinate (P) at (35:\R);
					\filldraw[blue!60!cyan] (P) circle (2.2pt);
					\draw[dashed] (C) -- (P) node[midway, below] {$r$};
					
					% velocidade tangencial
					\draw[->, thick, magenta!85!black]
					(P) -- +({-0.95*sin(35)},{0.95*cos(35)})
					node[above] {$\vec v$};
					
					% aceleração radial
					\draw[->, thick, orange!90!black]
					(P) -- ($(P)!0.5!(C)$)
					node[midway, below] {$\vec a_r$};
					
					% aceleração tangencial
					\draw[->, thick, green!60!black]
					(P) -- +({-0.55*sin(35)},{0.55*cos(35)})
					node[right] {$\vec a_t$};
					
					% sentido angular
					\draw[->, thick]
					(2.0,0.15) arc[start angle=5,end angle=55,radius=1.95];
					\node at (2.22,1.0) {$\omega$};
					
					\node[align=center, font=\scriptsize] at (0,-2.35)
					{Duas componentes\\ da aceleração.};
				\end{tikzpicture}
			\end{column}
		\end{columns}
		
	\end{frame}
	
	
	% SLIDE 13
	\begin{frame}{Interpretação vetorial de \(\omega\) e \(\alpha\)}
		\scriptsize
		
		Até agora usamos \(\omega\) e \(\alpha\) como grandezas escalares com sinal.
		Mas, em geral, a \textbf{velocidade angular} e a \textbf{aceleração angular}
		podem ser representadas por vetores:
		$
		\vec\omega, \qquad \vec\alpha.
		$
		
		\begin{columns}[T,onlytextwidth]
			
			\begin{column}{0.56\textwidth}
				
				\begin{itemize}
					\item \(\vec\omega\) não aponta na direção do movimento da partícula.\(\vec\omega\) aponta ao longo do \textbf{eixo de rotação};
					\item o sentido de \(\vec\omega\) é dado pela \textbf{regra da mão direita};
					\item \(\vec\alpha\) indica como \(\vec\omega\) varia no tempo.
				\end{itemize}
				
				\vspace{0.1cm}
				
				Para uma rotação no plano \(xy\), o eixo de rotação é o eixo \(z\):
				\[
				\omega>0
				\quad \Rightarrow \quad
				\vec\omega \text{ aponta no sentido positivo de } z,
				\]
				\[
				\omega<0
				\quad \Rightarrow \quad
				\vec\omega \text{ aponta no sentido negativo de } z.
				\]
				
				\vspace{0.1cm}
				
				\textcolor{red}{\textbf{No eixo fixo:}} podemos continuar a usar apenas
				\(\omega\) e \(\alpha\), com sinal.
			\end{column}
			
			\begin{column}{0.44\textwidth}
				\centering
				\begin{tikzpicture}[scale=0.85,>=stealth]
					
					%========================
					% Caso anti-horario
					%========================
					\begin{scope}[shift={(0,1.25)}]
						\coordinate (C) at (0,0);
						\def\R{1.05}
						
						\draw[thick, orange!60!brown] (C) circle (\R);
						\filldraw[ponto] (C);
						
						% sentido anti-horario
						\draw[->, thick, orange!80!black]
						(15:\R) arc[start angle=15,end angle=90,radius=\R];
						
						% eixo z saindo do plano
						\node[blue!80!black, font=\Large] at (C) {\(\odot\)};
						\node[right=0.9cm, blue!80!black] at (C) {\(\vec\omega\)};
						
						\node[font=\scriptsize] at (0,-1.45)
						{anti-horário: \(\omega>0\)};
						
						\node[font=\scriptsize, blue!80!black] at (0,1.45)
						{\(\vec\omega\) sai do plano};
					\end{scope}
					
					%========================
					% Caso horario
					%========================
					\begin{scope}[shift={(0,-2.55)}]
						\coordinate (C) at (0,0);
						\def\R{1.05}
						
						\draw[thick, orange!60!brown] (C) circle (\R);
						\filldraw[ponto] (C);
						
						% sentido horario
						\draw[->, thick, orange!80!black]
						(90:\R) arc[start angle=90,end angle=15,radius=\R];
						
						% eixo z entrando no plano
						\node[red!80!black, font=\Large] at (C) {\(\otimes\)};
						\node[right=0.9cm, red!80!black] at (C) {\(\vec\omega\)};
						
						\node[font=\scriptsize] at (0,-1.45)
						{horário: \(\omega<0\)};
						
						\node[font=\scriptsize, red!80!black] at (0,1.35)
						{\(\vec\omega\) entra no plano};
					\end{scope}
					
				\end{tikzpicture}
			\end{column}
		\end{columns}
		
		
		
	\end{frame}
	
	% SLIDE FINAL
	\begin{frame}{Resumo: da descrição linear à descrição angular}
		\scriptsize
		
		\begin{columns}[T]
			\begin{column}{0.5\textwidth}
				\textcolor{blue}{\textbf{Correspondências}}
				\[
				s \leftrightarrow \theta,
				\qquad
				v \leftrightarrow \omega,
				\qquad
				a_t \leftrightarrow \alpha
				\]
				
				\begin{align*}
					s &= r\theta, \qquad v = r\omega,\\
					a_t &= \alpha r, \qquad
					a_r = \frac{v^2}{r}=\omega^2 r.
				\end{align*}
				
				\vspace{0.1cm}
				\textcolor{blue}{\textbf{MCU}}
				\begin{align*}
					\omega &= \text{constante}, \qquad \alpha = 0,\\
					T &= \frac{2\pi}{\omega}=\frac{2\pi r}{v},\,a_c = \frac{v^2}{r}=\omega^2 r.
				\end{align*}
			\end{column}
			
			\begin{column}{0.5\textwidth}
				\textcolor{blue}{\textbf{MCUV}}
				\begin{align*}
					\omega(t) &= \omega(0)+\alpha t,\\
					\theta(t) &= \theta(0)+\omega(0)t+\frac{1}{2}\alpha t^2,\\
					\omega^2(t) &= \omega^2(0)+2\alpha\bigl(\theta(t)-\theta(0)\bigr).
				\end{align*}
				
				\vspace{0.1cm}
				\textcolor{blue}{\textbf{Dinâmica radial}}
				\begin{equation*}
					\sum F_r = m a_r = m\frac{v^2}{r}=m\omega^2 r.
				\end{equation*}
				
				\vspace{0.1cm}
				\textcolor{red}{\textbf{Ideia final:}}
				\begin{itemize}
					\item \(a_t\) altera a rapidez,\,\(a_r\) altera a direção;
					\item no MCU só há aceleração radial;
					\item no MCUV há aceleração radial e tangencial;
				\end{itemize}
			\end{column}
		\end{columns}
		
		
		
	\end{frame}
	
	
\end{document}